%
% 1-zahlen.tex -- 1. Komplexe Zahlen
%
% (c) 2025 Prof Dr Andreas Müller
%
\section{Komplexe Zahlen
\label{buch:kmoplex:section:zahlen}}
\kopfrechts{Komplexe Zahlen}
Die Lösungsformel für die quadratische Gleichung
\begin{equation}
x^2 + q = px
\qquad\text{oder}\qquad
x^2 - px + q = 0
\label{buch:komplex:zahlen:eqn:qeq}
\end{equation}
\index{quadratische Gleichung}%
war bereits in babylonischer Zeit etwa 2000 BC bekannt.
Auf Keilschrifttafeln aus der dritten Dynastie von Ur findet man den
\index{Keilschrifttafel}%
folgenden Algorithus:
\begin{enumerate}
\item Berechne die Hälfte von $p$: $\displaystyle\frac{p}{2}$.
\item Quadriere das Resultat: $\displaystyle\biggl(\frac{p}{2}\biggr)^2$.
\item Subtrahiere $q$: $\displaystyle\biggl(\frac{p}2\biggr)^2-q$.
\item Finde die (positive) Quadratwurzel mithilfe einer Tabelle von
Quadratzahlen:
\[
\!\sqrt{\rlap{\phantom{\bigg|}}\smash{\biggl(\frac{p}2\biggr)^2-q}}.
\]
\item Addiere die Resultate der Schritte 1 und 4:
$\displaystyle\frac{p}2+\!\sqrt{\rlap{\phantom{\bigg|}}\smash{\biggl(\frac{p}2\biggr)^2-q}}$.
\end{enumerate}
Der Algorithmus trägt sowohl in der
Formulierung~\eqref{buch:komplex:zahlen:eqn:qeq}
wie auch in der Durchführung der Tatsache Rechnung, dass die
Babylonier die negativen Zahlen noch nicht kannten.
Vor allem im Schritt 4 wird deutlich, dass nur positive Quadratwurzeln
in Betracht gezogen wurden.
Dies hat zur Folge, dass der Algorithmus nur eine Lösung finden kann,
selbst wenn wie im Falle der Gleichung
\begin{equation}
x^2 + 3 = 4x
\qquad
\Rightarrow
\qquad
p=4,\;q=3
\label{buch:komplex:zahlen:eqn:qepos}
\end{equation}
zwei positive, ganzzahlige Lösungen existieren.
Im Falle der Gleichung~\eqref{buch:komplex:zahlen:eqn:qepos}
ergibt der Algorithmus
\begin{align*}
x
&=
\frac{p}2
+
\!\sqrt{\rlap{$\phantom{\bigg|}$}\smash{\biggl(\frac{p}{2}\biggr)^2-q}}
=
\frac{4}{2}
+
\!\sqrt{\rlap{$\phantom{\biggl|}$}\smash{\biggl(\frac{4}{2}\biggr)^2-3}}
=
2
+
\!\sqrt{2^2-3}
=
2
+
\!\sqrt{1}
=
3.
\end{align*}
Die moderne Lösungsformel liefert dagegen die zwei Lösungen
\index{Losungsformel@Lösungsformel}%
\begin{align*}
x_\pm
&=
\frac{p}2
\pm
\!\sqrt{\rlap{$\phantom{\bigg|}$}\smash{\biggl(\frac{p}{2}\biggr)^2-q}}
=
\frac{4}{2}
\pm
\!\sqrt{\rlap{$\phantom{\biggl|}$}\smash{\biggl(\frac{4}{2}\biggr)^2-3}}
=
2\pm1
=
\begin{cases}
3\\
1.
\end{cases}
\end{align*}
Beide Lösung erfüllen ganz offensichtlich die ursprüngliche
Gleichung
\begin{align*}
x_+^2 + 3 &= 4x_+
&&\Leftrightarrow&
3^2 + 3 &= 4\cdot 3
&&\Leftrightarrow&
12&=12,
\\
x_-^2 + 3 &= 4x_-
&&\Leftrightarrow&
1^2 + 3 &= 4\cdot 1
&&\Leftrightarrow&
4&=4,
\end{align*}
auch wenn die zweite Lösung nicht vom Algorithmus gefunden werden kann.

In Schritt~4 wird verlangt, dass die Quadratwurzel mithilfe einer
Tabelle von Quadratzahlen gefunden werden muss.
\index{Quadratwurzel}%
Es können also nur Lösungen gefunden werden, für die die Quadratwurzel
im Schritt~4 eine in babylonischer Zeit bekannte Zahl ergab, also
\index{Babylon}%
im kompliziertesten Fall ein Bruch oder, wie wir heute sagen würden,
eine rationale Zahl.
Die Gleichung
\begin{align}
x^2 + 4x &= 2
&&\Rightarrow&
p&=4,\; q=2
&&\Rightarrow&
x
&=
\frac{p}{2}
+
\!\sqrt{\rlap{$\phantom{\bigg|}$}\smash{\biggl(\frac{p}2\biggr)^2-q}}
\label{buch:komplex:zahlen:eqn:qeunloesbar}
\\
&&&&&&&&
&=
\frac{4}{2}
+
\!\sqrt{\rlap{$\phantom{\bigg|}$}\smash{\biggl(\frac{4}2\biggr)^2-2}}
\notag
\\
&&&&&&&&
&=
2
+
\!\sqrt{2^2-2}
=
2
+
\!\sqrt{2}
\notag
\end{align}
Die Zahl $2$ kommt in keiner Tabelle von Quadratzahlen vor.
\index{Quadratzahl}%
Die Gleichung~\eqref{buch:komplex:zahlen:eqn:qeunloesbar} ist daher für die
babylonische Mathematik unlösbar.
Die moderne Mathematik muss also das Problem lösen, den zur Verfügung
stehenden Zahlenbereich so zu erweitern, dass Gleichungen
wie \eqref{buch:komplex:zahlen:eqn:qeunloesbar} lösbar sind.

Die genauere Inspektion der ursprünglichen
Gleichung~\eqref{buch:komplex:zahlen:eqn:qeq} fördert noch
weitere Schwierigkeiten des Lösungsalgorithmus zu Tage, zum Beispiel
den Fall $p=0$, $q=1$, der auf die Lösung $x=\sqrt{-1}$ führt.
in diesem Abschnitt sollen daher Techniken entwickelt werden, wie die
rationalen Zahlen so erweitert werden können, dass beliebige polynomiale
Gleichungen gelöst werden können.

%
% 11-reell.tex -- Erweiterung von Q durch die reellen Zahlen
%
% (c) 2026 Prof Dr Andreas Müller
%
\subsection{Erweiterung von $\mathbb{Q}$ durch die reellen Zahlen}
Schon im Altertum war bekannt, dass die Quadratwurzel $\!\sqrt{2}$ von 
$2$ nicht rational ist.
Diese Entdeckung, die normalerweise den Pythatgoräern zugeschrieben wird,
\index{Pythatgoraer@Pythagoräer}%
widersprach der verbreiteten Ansicht, dass sich alle Zahlen durch
Brüche ausdrücken lassen.
In diesem Abschnitt werden zwei Konstruktionen betrachtet, mit denen
sich eine Erweiterung von $\mathbb{Q}$ finden lässt, die $\!\sqrt{2}$
enthält.

%
% Die reellen Zahlen
%
\subsubsection{Die reellen Zahlen}
Die Quadratwurzel $\!\sqrt{2}$ ist bei weitem nicht die einzige irrationale
Zahl.
Jede ganze Zahl, deren Primfaktorzerlegung einen Primfaktor in ungerader
Anzahl enthält, hat eine irrationale Quadratwurzel.
Auch viele kubische Wurzeln, die eulersche Zahl $e$ oder die
Kreiszahl $\pi$ sind irrational, es lassen sich aber immer gute
Approximationen durch Brüche geben.
\index{Kreiszahl}%
\index{pi@$\pi$}%
Leonhard Euler hat beretis 1737 bewiesen, dass $e$ irrational ist.
\index{Euler, Leonhard}%
Für $\pi$ ist dies erst Johann Heinrich Lambert 1761 gelungen.
\index{Lambert, Johann Heinrich}%
Diese Beispiele deuten darauf hin, dass es eine grosse Menge von
Zahlen gibt, die beliebig genau durch Brüche approximiert werden
können, aber trotzdem keine Brüche sind.

Das Problem kann gelöst werden, indem man eine Menge $\mathbb{R}$ von
Zahlen definiert, die beliebig genau durch Brüche approximiert werden
können.
Dabei muss aber auch berücksichtigt werden, dass eine solche Zahl
verschiedene Approximationen haben kann.
Zum Beispiel kann $\pi$ sowohl durch Dezimalbrüche
\[
3
, 3.1
, 3.14
, 3.141
, 3.1415
, 3.14159
, 3.141592
, 3.1415926
, 3.14159265
%, 3.141592659
%, 3.1415926592
%, 3.14159265926
\to
\pi
\]
als auch durch die Kettenbruchentwicklung
\index{Kettenbruch}%
\[
\frac{3}{1}
\approx
3.{\color{gray}00},
\frac{22}{7} \approx 3.14{\color{gray}28},
\frac{333}{106}\approx 3.1415{\color{gray}09},
\frac{355}{113}\approx 3.141592{\color{gray}92}
\to
\pi
\]
approximiert werden.
Es reicht daher nicht, eine Approximation zu haben.
Es wird ausserdem ein Kriterium benötigt, mit welchem man entscheiden
kann, ob zwei Approximationen die gleiche Zahl approximieren.

\begin{definition}[Cauchy-Folge]
Eine Folge $(a_n)_{n\in\mathbb{N}}$ von Zahlen $a_n\in\mathbb{Q}$ heisst
\emph{Cauchy-Folge}, wenn für jedes $\varepsilon>0$ ein $N\in\mathbb{N}$
\index{Cauchy-Folge}%
gibt derart, dass
\[
|a_n-a_m| < \varepsilon
\]
für alle $m,n> N$ gilt.
\end{definition}

Die Dezimalbruchapproximation von $\pi$ ist eine Cauchyfolge.
Sei $a_n$ der endliche Dezimalbruch mit $n$ Nachkommastellen, die
mit der Dezimalbruchentwicklung übereinstimmen.
Dann ist
$a_n-a_m$ ein Dezimalbruch, dessen erste $\min(n,m)$ Nachkommastellen
$0$ sind.
Insbesondere ist
\[
|a_n-a_m| < 10^{-\min(n,m)}.
\]
Sei also $\varepsilon>0$ gegeben und $N\in\mathbb{N}$ derart,
dass $10^{-N}<\varepsilon$.
Falls $n,m>N$ sind, folgt
\[
|a_n-a_m|
<
10^{-\min(n,m)}
<
10^{-N}
<
\varepsilon.
\]
Damit ist gezeigt, dass $(a_n)_{n\in\mathbb{N}}$ eine Cauchy-Folge ist.

Endliche und unendliche Dezimalbrüche konvergieren gegen rationale Zahlen
im Sinne der folgenden Definition.

\begin{definition}[konvergente Folge]
\index{konvergente Folge}%
Eine Cauchy-Folge $(a_n)_{n\in\mathbb{N}}$ konvergiert gegen $a\in\mathbb{Q}$,
auch geschrieben als
\[
\lim_{n\to\infty} a_n = a
\qquad\text{oder}\qquad
a_n\to a,
\]
wenn für jedes $\varepsilon>0$ ein $N\in\mathbb{N}$ existiert derart, dass
\[
|a_n-a|<\varepsilon
\]
für $n>N$.
\end{definition}

Nicht jede Cauchy-Folge konvergiert gegen eine rationale Zahl.
Die Approximationsfolge von $\!\sqrt{2}$ ist eine Cauchy-Folge, die nicht
gegen eine rationale Zahl konvergiert.

Mit Cauchy-Folgen kann man auch Rechnen.
In der Analysis wird der folgende Satz bewiesen.

\begin{satz}
Seien $(a_n)_{n\in\mathbb{N}}$ und $(b_n)_{n\in\mathbb{N}}$ Cauchy-Folgen.
Dann sind $(a_n+b_n)_{n\in\mathbb{N}}$, $(a_n-b_n)_{n\in\mathbb{N}}$ 
und $(a_nb_n)_{n\in\mathbb{N}}$ Cauchy-Folgen.
Falls $a_n\to a$ und $b_n\to b$, dann gilt auch
$a_n\pm b_n\to a\pm b$ und $a_nb_n\to ab$.
Falls $b\ne 0$ ist, ist auch $(a_n/b_n)_{n\in\mathbb{N}}$ ein Cauchy-Folge
und $a_n/b_n \to a/b$.
\end{satz}

\begin{definition}[Nullfolge]
Eine Cauchy-Folge $(a_n)_{n\in\mathbb{N}}$ heisst \emph{Nullfolge}, wenn
sie gegen $0$ konvergiert.
\index{Nullfolge}%
\end{definition}

Die Rechenoperationen und der Begriff einer Nullfolge ermöglichen jetzt,
Cauchy-Folgen miteinander zu identifizieren, die den gleichen Grenzwert
haben, selbst wenn dieser nicht eine rationale Zahl ist.

\begin{definition}[Äquivalenz von Cauchy-Folgen]
Zwei Cauchy-Folgen
$(a_n)_{n\in\mathbb{N}}$
und
$(b_n)_{n\in\mathbb{N}}$
heissen \emph{äquivalent}, in Zeichen $(a_n)\sim (b_n)$ wenn
$(a_n-b_n)_{n\in\mathbb{N}}$ eine Nullfolge ist.
\index{aquivalente@äquivalente Cauchy-Folgen}%
\end{definition}

Wenn zwei Cauchy-Folgen äquivalent sind, dann approximieren sie 
die gleiche Zahl.
Die gesuchte Erweiterung der Menge der rationalen Zahlen $\mathbb{Q}$
besteht also aus Cauchy-Folgen von rationalen Zahlen, wobei jedoch
äquivalente Folgen die gleiche Zahl repräsentieren.

\begin{definition}[Äquivalenzklassen]
\index{Aquivalenzklasse@Äquivalenzklasse}%
Eine \emph{Äquivalenzklasse} $X$ der Äquivalenzrelation $\sim$ von
\index{Aquivalenzrelation@Äquivalenzrelation}%
Cauchy-Folgen ist eine Menge von äquivalenten Cauchy-Folgen.
\end{definition}

Zwei Cauchy-Folgen
$(a_n)_{n\in\mathbb{N}}\in X$ 
und
$(b_n)_{n\in\mathbb{N}}\in X$ 
in einer Äquivalenzklasse $X$ von rationalen Cauchy-Folgen unterscheiden
sich nur durch eine Nullfolge, es gilt also
\[
\lim_{n\to\infty} (a_n-b_n) = 0.
\]
Die reellen Zahlen können jetzt als Äquivalenzklassen von Cauchy-Folgen
definiert werden.

\begin{definition}[Reelle Zahlen]
Die Menge der reellen Zahlen $\mathbb{R}$ ist die Menge der Äquivalenzklassen
von rationalen Cauchy-Folgen bezüglich der Äquivalenzrelation $\sim$.
\end{definition}

Diese Konstruktion zeigt, dass man mit reellen Zahlen rechnen kann, indem
man geeignete Approximationen dafür verwendet.
Welche Approximation verwendet wird, ist egal.
Wenn die Approximationsfolgen äquivalent sind, unterscheiden sich auch
Resultate von Rechnungen mit den Folgen durch Nullfolgen.
Die reellen Zahlen enthalten nach Konstruktionen alle Quadratwurzeln
und die Zahlen $e$ und $\pi$.
Jede Gleichung, für deren Lösung eine Approximationsfolge existiert,
kann daher in $\mathbb{R}$ auch gelöst werden.
So kann man zum Beispiel die Lösung der Gleichung
\[
y = 1 + \frac1y
\]
dadurch approximieren, dass man die Folge $(y_n)_{n\in\mathbb{N}}$
durch
\[
y_0 = 1,
\quad
y_{n+1} = 1 + \frac1{y_n}
\]
bildet.
$(y_n)_{n\in\mathbb{N}}$ ist eine Cauchy-Folge, die gegen das Verhältnis
$\varphi=(\!\sqrt{5}+1)/2$ des goldenen Schnittes konvergiert (siehe
Tabelle~\ref{buch:komplex:zahlen:table:gold}, zweite Spalte).
%
% table-gold.tex
%
% (c) 2026 Prof Dr Andreas Müller
%
\begin{table}
\def\p{\phantom{-}}
\centering
\begin{tabular}{|>{$}r<{$}|>{$}l<{$}|>{$}l<{$}|}
\hline
&\begin{minipage}[t]{3.5cm}\raggedright
\setlength{\abovedisplayskip}{0pt}
\setlength{\belowdisplayskip}{0pt}
\raisebox{3pt}{\strut}Rekursive Approxima- tion
\[
 y_{n+1}=1+\frac{1}{y_n}
\]
für $\varphi$
\end{minipage}&\begin{minipage}[t]{3.5cm}\raggedright
Newton-Approximation für Nullstellen von $f(x)=x^2+1$
\end{minipage}
\\[44pt]
\hline
n & y_n\raisebox{3pt}{\mathstrut} & x_n \\[2pt]
\hline
 1 & 2                 & \p0.508034707497476 \raisebox{4pt}{\mathstrut}
\\
 2 & 1.500000000000000             & -0.730167373438399 \\
 3 & 1.\underline{6}66666666666667 &\p0.319690815934174 \\
 4 & 1.\underline{6}00000000000000 & -1.404165739925735 \\
 5 & 1.\underline{6}25000000000000 & -0.345999548896763 \\
 6 & 1.\underline{61}5384615384615 &\p1.272088814812140 \\
 7 & 1.\underline{61}9047619047619 &\p0.242990090617789 \\
 8 & 1.\underline{61}7647058823529 & -1.936202034966178 \\
 9 & 1.\underline{618}181818181818 & -0.709863503540628 \\
10 & 1.\underline{61}7977528089888 &\p0.349429012666962 \\
11 & 1.\underline{6180}55555555556 & -1.256191291052456 \\
12 & 1.\underline{6180}25751072961 & -0.230067094013908 \\
13 & 1.\underline{61803}7135278515 &\p2.058245522488201 \\
14 & 1.\underline{61803}2786885246 &\p0.786197418015100 \\
15 & 1.\underline{61803}4447821682 & -0.242873870579829 \\
16 & 1.\underline{618033}813400125 &\p1.937244794475075 \\
17 & 1.\underline{61803}4055727554 &\p0.710523884635529 \\
18 & 1.\underline{6180339}63166706 & -0.348444169203709 \\
19 & 1.\underline{6180339}98521804 &\p1.260728028475482 \\
20 & 1.\underline{61803398}5017358 &\p0.233767770871423 \\
21 & 1.\underline{6180339}90175597 & -2.021990939507587 \\
22 & 1.\underline{618033988}205325 & -0.763714440828034 \\
23 & 1.\underline{618033988}957902 &\p0.272837745754085 \\
24 & 1.\underline{618033988}670443 & -1.696172136911838 \\
25 & 1.\underline{6180339887}80243 & -0.553304666781480 \\
26 & 1.\underline{6180339887}38303 &\p0.627008940439556 \\
27 & 1.\underline{6180339887}54322 & -0.483932324938967 \\
28 & 1.\underline{61803398874}8204 &\p0.791236155774208 \\
29 & 1.\underline{6180339887}50541 & -0.236304511027909 \\
30 & 1.\underline{618033988749}648 &\p1.997761646531475 \\[2pt]
\hline
\infty & 1.618033988749895 & \raisebox{4pt}{\mathstrut}\\[2pt]
\hline
\end{tabular}
\caption{Approximation der Zahl $\varphi=(1+\!\sqrt{5})/2$
durch die Folge $(y_n)_{n\in\mathbb{N}}$, die durch $y_0=1$ und
$y_{n+1}=1+1/y_n$ rekursiv definiert ist (Spalte $y_n$).
In der Spalte $x_n$ sind die Newton-Approximationen für die
Funktion $f(x) = x^2 + 1$, rekursiv definiert durch
$x_{n+1} = x_n -  f(x_n)/f'(x_n)$, dargestellt, die keinerlei
Anzeichen von Konvergenz zeigen.
\label{buch:komplex:zahlen:table:gold}}
\end{table}


Approximationsfolgen für die Nullstellen einer Funktion $f(x)$ können
zum Beispiel durch das Newton-Verfahren gefunden werden.
\index{Newton-Verfahren}%
\index{Nullstelle}%
Dazu wählt man eine Anfangsschätzung $x_0$ und bildet dann die 
Folge
\[
x_{n+1}
=
x_n - \frac{f(x_n)}{f'(x_n)}.
\]
Falls $x_0$ genügend nahe bei einer Nullstelle $\hat{x}$ liegt, wird die
Folge $(x_n)_{n\in\mathbb{N}}$ gegen die Nullstelle konvergieren:
\[
\lim_{n\to\infty} x_n
=
\hat{x}.
\]
Wendet man dies auf die Funktion $f(x)=x^2+1$ an, welche $f(x) \ge 1$
für alle $x\in\mathbb{R}$ keine reelle Nullstelle hat, erhält man eine
Folge, die keine Cauchy-Folge ist
(Abbildung~\ref{buch:komplex:zahlen:table:gold}, dritte Spalte).
Die Gleichung
\begin{equation}
x^2 + 1 = 0
\qquad\Leftrightarrow\qquad
x^2 = -1
\end{equation}
hat also keine Lösung, die mit einer Approximationsfolge gefunden werden
kann.


%
% 12-algebraisch.tex -- Algebraische Erweiterungen von Q
%
% (c) 2026 Prof Dr Andreas Müller
%
\subsection{Algebraische Erweiterungen von $\mathbb{Q}$}
Die Vervollständigung der der rationalen Zahlen zu den reellen
Zahlen stellt sicher, dass jede Cauchy-Folge einen Grenzwert hat.
Sie führt aber zu einem neuen Problem, welches ein Hindernis für
Computeralgebrasysteme wird.
\index{Computeralgebra}%
Da reelle Zahlen nur durch ihre Approximationen definiert sind,
lässt sich damit grundsätzlich nur dann exakt rechnen, wenn auch
noch eine alternative Charakterisierung zur Verfügung steht.
Mit der Kreiszahl $\pi$ kann man nur desshalb exakt rechnen, weil
sie die kleinste positive Nullstelle der Sinusfunktion ist.

Für die Quadratwurzel $\!\sqrt{2}$, die früher als erstes Beispiel
\index{Quadratwurzel}%
einer irrationalen Zahl in den Vordergrund gestellt wurde, gibt
es eine solche alternative Charakterisierung.
Beim algebraischen Rechnen mit $\!\sqrt{2}$ verwendet man fast
instinktiv die Beziehung $(\!\sqrt{2})^2=2$, zum Beispiel
beim Ausmultiplizieren von
\[
(\!\sqrt{2}+1)^2
=
(\!\sqrt{2})^2
+2\!\sqrt{2}\cdot 1 
+ 1^2
=
2 + 2\!\sqrt{2} + 1
=
3+2\!\sqrt{2}.
\]

Wir wollen diese Beobachtung noch etwas expliziter machen und
bezeichnen mit $\alpha$ ein Objekt, für welches die Regeln
der Algebra ebenso gelten wie die Identität $\alpha^2-2=0$.
Durch die Ersetzung $\alpha^2 = 2$ kann jede Potenz von $\alpha$
auf entweder ein Vielfaches von $\alpha$ oder auf eine ganze Zahl
reduziert werden, es gilt nämlich
\[
\alpha^n
=
\begin{cases}
2^{\frac{n}2}         &\qquad\text{$n$ gerade}\\
2^{\frac{n-1}2}\alpha &\qquad\text{$n$ ungerade.}
\end{cases}
\]
Für Linearkombinationen $a+b\alpha$ und $c+d\alpha$ lassen sich
jetzt einfache Rechenregeln aufstellen:
\begin{align*}
(a+b\alpha)\pm(c+d\alpha) &= (a\pm b) + (c\pm d)\alpha
\\
(a+b\alpha)(c+d\alpha) &= (ac + 2bd) + (ad+bc)\alpha.
\end{align*}
Als Spezialfall der Produktregel findet man
\[
(a+b\alpha)(a-b\alpha)
=
a^2 - 2 b^2,
\]
daher heisst $a-b\alpha$ auch das zu $a+b\alpha$ \emph{konjugierte}
Element.
\index{konjugiert}%
Durch Erweitern mit $(c-d\alpha)$ kann man sogar den Kehrwert
\begin{align*}
\frac{1}{c+d\alpha}
&=
\frac{1}{c+d\alpha}
\cdot
\frac{c-d\alpha}{c-d\alpha}
=
\frac{c-d\alpha}{c^2-2d^2}
=
\frac{c}{c^2-2d^2} - \frac{d}{c^2-2d^2}\alpha
\intertext{und damit auch den Quotienten}
\frac{a+b\alpha}{c+d\alpha}
&=
\frac{ac-2bd}{c^2-2d^2} - \frac{ad-bc}{c^2-2d^2}\alpha
\end{align*}
definieren.

Alle vier Grundoperationen lassen sich also für alle Zahlen der Form 
$a+b\alpha$ ausschliesslich durch Verwendung der arithmetischen Gesetze
und der Identität $\alpha^2-2$ definieren und sie ergeben immer wieder
Zahlen dieser Form.
Es entsteht somit eine neue Menge
\[
\mathbb{Q}(\alpha)
=
\{a+b\alpha\mid a,b\in\mathbb{Q}\},
\]
in der die gleichen Rechengesetze gelten wie in $\mathbb{Q}$.
Die Menge $\mathbb{Q}$ ist eine echte Teilmenge, $\mathbb{Q}(\alpha)$
ist also eine echte Erweiterung von $\mathbb{Q}$.
In $\mathbb{Q}(\alpha)$ lässt sich jetzt die Gleichung $X^2-2=0$ lösen.
Es ist nämlich
\[
(X+\alpha)(X-\alpha)
=
X^2-\alpha^2
=
X^2 - 2
=
0,
\]
die beiden Lösungen der Gleichung $X^2-2=0$ sind also $\pm\alpha$.

Die Konstruktion lässt sich noch etwas verallgemeinern.
Ein \emph{normiertes Polynom} ist ein Polynom dessen führender Koeffizient
\index{normiertes Polynom}%
\index{Polynom!normiert}%
$1$ ist.
Verwendet man statt des Polynoms $X^2-2$ ein beliebiges normiertes
Polynom $m(X)$ aus der Menge
\[
\mathbb{Q}[X]
=
\{
X^n+m_{n-1}X^{n-1}+\dots+m_1X+m_0 
\mid
m_0,\dots,m_{n-1}\in\mathbb{Q}
\}
\]
der Polynome mit rationalen Koeffizienten ersetzt wird\footnote{Im Folgenden
werden wir die Variable eines Polynoms mit grossen Buchstaben bezeichnen.
Sie wird damit klar von den Werten unterschieden, die man anstelle der
Variable ein setzen kann.
Die Variable $X$ kann man auch als ein transzendentes Element über dem
Koeffizientenkörper betrachten.
}.
Wieder versucht man ein neues Objekt $\alpha$ durch die Eigenschaft
$m(\alpha)=0$ zu charakterisieren.
Dies führt jedoch auf Schwierigkeiten, wenn sich $m$ als Produkt
$m=ab$ von Polynomen $a(X),b(X)\in\mathbb{Q}[X]$ schreiben lässt,
denn die beiden Faktoren $a(X)$ und $b(X)$ könnten Nullstellen mit
widersprüchlichen Eigenschaften definieren.

\begin{definition}[irreduzibles Polynom]
Das Polynom $m(X)\in\mathbb{Q}[X]$ heisst \emph{irreduzibel}, wenn es keine
\index{irreduzibel}%
Polynome vom Grad $>0$ gibt, die $m$ als Produkt haben.
\end{definition}

Man beachte, dass die Eigenschaft, irreduzibel zu sein, vom
Koeffizientenkörper abhängt.
Das Polynom $X^2-2$ ist irreduzibel in $\mathbb{Q}[X]$, es lässt sich
nicht in zwei Faktoren zerlegen.
Die einzig mögliche Faktorisierung wäre $X^2-2=(X+\!\sqrt{2})(X-\!\sqrt{2})$,
aber die Faktoren $X+\!\sqrt{2}$ und $X-\!\sqrt{2}$ haben nicht
alle Koeffizienten in $\mathbb{Q}$.
Das Polynom $X^2-2\in\mathbb{R}[X]$ ist also nicht irreduzibel als Polynom
mit reellen Koeffizienten.

In einem irreduziblen Polynom $m(X)$ kann der Koeffizient $m_0$ nicht
verschwinden, denn dann liesse sich
\[
m
=
(X^{n-1} + m_{n-1}X^{n-2}+\dots + m_1)X
\]
als Produkt zweier Polynome vom Grad $n-1$ bzw.~1 schreiben.

Mit einem irreduziblen Polynom $m(X)$ definieren wir jetzt wieder das
Objekt $\alpha$ mit der Eigenschaft $m(\alpha)=0$.
$m$ heisst auch das \emph{Minimalpolynom} von $\alpha$.
\index{Minimalpolynom}%
Ausgeschrieben bedeutet dies
\[
m(\alpha)
=
\alpha^n
+
m_{n-1}\alpha^{n-1}
+
\dots
+
m_1\alpha
+
m_0
=
0
\quad\Rightarrow\quad
\alpha^n
=
-m_{n-1}\alpha^{n-1}
-
\dots
-
m_1\alpha
-
m_0.
\]
Man kann also jede Potenz grösser oder gleich $n$ von $\alpha$ durch
eine Linearkombination von Potenzen $\alpha^k$ mit $0\le k< n$ 
ausdrücken.

\begin{definition}[algebraische Erweiterung]
Sei $m(X)\in\mathbb{Q}[X]$ ein irreduzibles Polynom vom Grad $n$
der Form
\[
m(X)
=
X^n + m_{n-1}X^{n-1} + \dots + m_1X + m_0.
\]
Die Menge
\begin{equation}
\mathbb{Q}(\alpha)
=
\{
a_0 + a_1\alpha + \dots + a_{n-1}\alpha^{n-1}
\mid
a_k\in\mathbb{Q}\text{ mit } 0\le k<n
\}
\label{buch:komplex:zahlen:def:eqn:Qa}
\end{equation}
heisst die \emph{algebraische Erweiterung} von $\mathbb{Q}$ um $\alpha$.
\index{Erweiterung, algebraisch}%
\index{algebraische Erweiterung}%
\end{definition}

Wieder sollen die Rechengesetze durch Übertragung der gewohnten
Regeln entstehen, jedoch müssen alle höheren Potenzen von $\alpha$
jeweils wieder auf Potenzen $\alpha^k$, $0\le k<n$, reduziert werden.
Da der Koeffizient $m_0$ des irreduziblen Polynoms $m$ nicht verschwindet,
kann man die Gleichung $m(\alpha) = 0$ mit $\alpha^{-1}$ multiplizieren
und nach $\alpha^{-1}$ auflösen:
\begin{align*}
m(\alpha)\cdot \alpha^{-1}
=
0
&=
\alpha^{n-1}
+
m_{n-1}\alpha^{n-2}
+
\dots
+
m_1
+
m_0\alpha^{-1}
\\
\alpha^{-1}
&=
-\frac{1}{m_0}
(
\alpha^{n-1}
+
m_{n-1}\alpha^{n-2}
+
\dots
+
m_1
)
\\
&=
-\frac{m_1}{m_0}
-\dots
-\frac{m_{n-1}}{m_0}\alpha^{n-2}
-\frac{1}{m_0}\alpha^{n-1}
\in
\mathbb{Q}(\alpha).
\end{align*}
Das reziproke Element ist also immer auch in $\mathbb{Q}(\alpha)$.
\index{reziprok}%

\begin{beispiel}
\label{buch:komplex:zahlen:bsp:kubikwurzel}
Die Gleichung $X^3-2=0$ lässt sich in $\mathbb{Q}$ nicht lösen.
Das Polynom $m(X)=X^3-2$ ist irreduzibel als Polynom in $\mathbb{Q}[X]$.
Die Erweiterung $\mathbb{Q}(\alpha)$ um das Objekt $\alpha$ mit
\index{Kubikwurzel}%
der Eigenschaft $\alpha^3-2=0$ ist also die Menge
\[
\mathbb{Q}(\alpha)
=
\{
a_0+a_1\alpha + a_2\alpha^2
\mid
a_0,a_1,a_2\in\mathbb{Q}
\}.
\]
Die Additionsregeln sind nicht weiter anspruchsvoll, wir berechnen
daher nur das Produkt von zwei Elementen $a,b\in\mathbb{Q}(\alpha)$:
\begin{align*}
ab
&=
(a_0+a_1\alpha+a_2\alpha^2)
(b_0+b_1\alpha+b_2\alpha^2)
\\
&=
a_0b_0
+
(a_0b_1+a_1b_0)\alpha
+
(a_0b_2+a_1b_1+a_2b_0)\alpha^2
+
(a_1b_2+a_2b_1){\color{darkred}\alpha^3}
+
a_2b_2{\color{darkred}\alpha^4}.
\intertext{Die Potenzen $\alpha^3$ und $\alpha^4$ müssen durch Anwendung der
Reduktionsregel $\alpha^3=2$ auf}
&=
a_0b_0
+
(a_0b_1+a_1b_0)\alpha
+
(a_0b_2+a_1b_1+a_2b_0)\alpha^2
+
(a_1b_2+a_2b_1){\color{darkred}\cdot 2}
+
a_2b_2{\color{darkred} \cdot 2\alpha}
\intertext{reduziert werden.
Das Produkt ist daher}
ab
&=
(a_0b_0 + {\color{darkred}2}a_1b_2+{\color{darkred}2}a_2b_1)
+
(a_0b_1+a_1b_0 + {\color{darkred}2}a_2b_2)\alpha
+
(a_0b_2+a_1b_1+a_2b_0)\alpha^2
\in\mathbb{Q}(\alpha).
\end{align*}
Das zu $\alpha$ reziproke Element ist
\[
\alpha^{-1}
=
\frac{1}{2}\alpha^2,
\]
tatsächlich ist
\[
\alpha\cdot \alpha^{-1}
=
\alpha
\cdot
\frac{1}{2}\alpha^2
=
\frac{\alpha^3}{2}
=
\frac{2}{2}
=
1.
\]
In $\mathbb{Q}(\alpha)$ kann man jetzt auch das Polynom $X^3-2$ 
faktorisieren, indem man durch den Faktor $(X-\alpha)$ dividiert.
\index{Polynomdivision}%
Die Polynomdivision ergibt
\[
\renewcommand{\arraycolsep}{1.0pt}
\begin{array}{rcrcrcrcrcrcrcrcr}
\llap{$($}
X^3 & &           & &            &-&        2\rlap{$)$}&\phantom{(}:\phantom{)}&\llap{$($} X &-& \alpha\rlap{$)$} &\phantom{(}=\mathstrut& X^2 &+& \alpha X &+& \alpha^2 \\
X^3 &-&\alpha X^2 & &            & &         & &   & &        & &     & &          & &          \\
\cline{1-7}
    & &\alpha X^2\raisebox{4pt}{\mathstrut} & &            &-&        2& &   & &        & &     & &          & &          \\
    & &\alpha X^2 &-& \alpha^2 X & &         & &   & &        & &     & &          & &          \\
\cline{3-7}
    & &           & & \alpha^2 X\raisebox{4pt}{\mathstrut} &-&        2& &   & &        & &     & &          & &          \\
    & &           & & \alpha^2 X &-& \alpha^3& &   & &        & &     & &          & &          \\
\cline{5-7}
    & &           & &            & & \alpha^3\raisebox{4pt}{\mathstrut}\rlap{$\mathstrut-2 = 0$.}& &   & &        & &     & &          & &
\end{array}
\]
Tatsächlich ergibt Ausmultiplizieren zur Kontrolle
\begin{align}
(X-\alpha)
(X^2 + \alpha X + \alpha^2)
&=
X^3 +\alpha X^2 +\alpha^2 X
-
\alpha
X^2
-
\alpha^2 X
-
\alpha^3
\label{buch:komplex:zahlen:eqn:faktorisierung}
\\
&=
X^3-\alpha^3
\notag
\\
&=
X^3-2.
\notag
\end{align}
In der üblichen algebraischen Notation würde man $\alpha=\sqrt[3]{2}$
für die Kubikwurzel schreiben.

Wir erwarten aber, dass die Gleichung $X^3-2=0$ zwei weitere Lösungen
hat, die sich aus der Faktorisierung
\eqref{buch:komplex:zahlen:eqn:faktorisierung}
berechnen lassen sollten.
Das Polynom $X^2 +\alpha X + \alpha^2$ ist allerdings irreduzibel
als Polynom in $\mathbb{Q}(\alpha)[X]$, es definiert also eine neue
Erweiterung um eine Zahl $\beta$, die das Minimalpolynom
$b(X)=X^2+\alpha X + \alpha^2$ hat.
Die Faktorisierung ist also erst in $\mathbb{Q}(\alpha)(\beta)
=\mathbb{Q}(\alpha,\beta)$ möglich.

Die konventionelle Schreibweise für die Nullstelle des Polynoms
$b$ ergibt sich aus der Lösungsformel für quadratische Gleichungen.
Man muss
\begin{equation}
\beta
=
-\frac{\alpha}{2}
+
\!\sqrt{
\rlap{$\phantom{\bigg|}$}
\smash{\biggl(\frac{\alpha}{2}\biggr)^{\raisebox{-1pt}{$\scriptstyle2$}}-\alpha^2}
}
=
-\frac{\alpha}{2}
+
\!\sqrt{-\frac34\alpha^2}
\label{buch:komplex:zahlen:eqn:cbrtnegativ}
\end{equation}
setzen.
\end{beispiel}

Das Beispiel zeigt, wie die Lösung einer Gleichung möglicherweise
das wiederholte Hinzufügen neuer Zahlen nötig macht.
Die Notation $\sqrt[3]{2}$ für das Element $\alpha$ ist gefährlich,
da sie bereits suggeriert, dass $\alpha$ eine reelle Zahl sein
soll.
Wir haben aber nur gefordert, dass $\alpha$ eine beliebige
Nullstelle des Polynoms sein soll.
Die Darstellung \eqref{buch:komplex:zahlen:eqn:cbrtnegativ} der
anderen Nullstellen zeigt für $\alpha=\sqrt[3]{2}$, dass diese
wegen
\[
-\frac34 \!\sqrt[3]{2} < 0
\]
keine reellen Zahlen sein können.
Dieses Argument verwendet die Ordnungsrelation in den reellen
Zahlen und insbesondere die Eigenschaft, dass negative Zahlen
keine reellen Quadratwurzeln haben.
\index{Ordnungsrelation}%
Auf die Ordnungsrelation wurde in der Konstruktion von
$\mathbb{Q}(\alpha)$ nirgends Bezug genommen.
Sogar die Konstruktion von $\mathbb{Q}(\!\sqrt{2})$ hätte genauso 
gut im $-\!\sqrt{2}$ an Stelle von $\alpha$ funktionieren können.
Mit rein arithmetischen Mitteln sind die beiden Nullstellen von
$X^2-2$ nicht unterscheidbar.

\begin{definition}[algebraisches Element]
Ein Element $\alpha$ heisst \emph{algebraisch über $\mathbb{Q}$},
wenn es ein Polynom $m(X)\in\mathbb{Q}[X]$ gibt, welches $m(\alpha)=0$
erfüllt.
\index{algebraisches Element}%
\index{Element, algebraisch}%
\end{definition}

Die Definition von $\sqrt{2}$ oder $\sqrt[3]{2}$ als
algebraische Elemente über $\mathbb{Q}$ ermöglicht also,
mit diesen Elementen exakt zu rechnen, ohne dass auf
Approximationen zurückgegriffen werden muss.
Die algebraische Definition ist daher die Methode der
Wahl für ein Computeralgebrasystem.

%
% Division im Erweiterungskörper
%
\subsubsection{Division in der algebraischen Erweiterung}
Die Konstruktion von $\alpha^{-1}$ in einer algebraischen Erweiterung
$\mathbb{Q}(\alpha)$ erklärt nicht, wie man durch ein allgemeines
Element wie in \eqref{buch:komplex:zahlen:def:eqn:Qa} dividiert.
Dazu muss man die multiplikative Inverse
\[
(a_0+a_1\alpha + \dots + a_{n-1}\alpha^{n-1})^{-1}
=
b_0+b_1\alpha + \dots + b_{n-1}\alpha^{n-1}
\]
bestimmen können.
Multiplizieren wir das Produkt
\[
(a_0+a_1\alpha + \dots + a_{n-1}\alpha^{n-1})
(b_0+b_1\alpha + \dots + b_{n-1}\alpha^{n-1})
=
1
\]
aus und ersetzen wir wiederholt Potenzen $\alpha^n$ durch
$-m_0-m_1\alpha-\dots-m_{n-1}\alpha^{n-1}$, entsteht ein
Gleichungssystem für die Koeffizienten $b_k$.
\index{lineares Gleichungssystem}%
Die Formulierung und Lösung dieses Gleichungssystems ist jedoch
in Matrizenschreibweise deutlich einfacher, wir verschieben daher
die allgemeine Berechnung der multiplikativen Inversen auf
Abschnitt~\ref{buch:komplex:subsection:matrixschreibweise}.

\begin{beispiel}
Wir betrachten die Körpererweiterung $\mathbb{Q}(\alpha)$, wobei
$\alpha$ das Minimalpolynom $m(X) = X^2+X+1$ hat.
Um die multiplikative Inverse von $a+b\alpha$ zu bestimmen,
suchen wir $x+y\alpha$ derart, dass
\begin{align*}
(a+b\alpha)(x+y\alpha)&=1
\\
ax + (bx+ay)\alpha + by\alpha^2 &= 1.
\intertext{Darin ersetzen wir $\alpha^2=-\alpha-1$ und erhalten}
ax + (bx+ay)\alpha - by(1+\alpha) &= 1 \\
(ax-by) + (bx+(a-b)y)\alpha &= 1
\end{align*}
Daraus lesen wir das Gleichungssystem
\index{lineares Gleichungssystem}%
\[
\left.
\renewcommand{\arraycolsep}{2pt}
\begin{array}{rcrcr}
ax &-&     by &=& 1 \\
bx &+& (a-b)y &=& 0
\end{array}
\quad\right\}
\qquad\text{mit der Lösung}\qquad
\begin{pmatrix}x\\y\end{pmatrix}
=
\frac{1}{b^2+ab-a^2}
\begin{pmatrix}
b-a\\b
\end{pmatrix}
\]
ab.
In der üblichen Schreibweise bedeutet dies, dass
\[
(a+b\alpha)^{-1}
=
\frac{(b-a)+b\alpha}{b^2+ab-a^2}
\]
die multiplikative Inverse ist.
\end{beispiel}

%
% Die Reihenfolge der Erweiterungen
%
\subsubsection{Die Reihenfolge der Erweiterungen}
Die Quadratwurzel $\!\sqrt{2}$ ist nicht in $\mathbb{Q}$, liegt aber in
der algebraischen Erweiterung $\mathbb{Q}(\!\sqrt{2})$, die mithilfe des
über $\mathbb{Q}$ irreduziblen Polynoms $m(X)=X^2-2\in\mathbb{Q}[X]$ 
konstruiert werden kann.

$\sqrt[4]{2}$ ist ein algebraisches Element über $\mathbb{Q}$, ist
aber selbst nicht Teil dieses Körpers.
Mithilfe des über $\mathbb{Q}$ irreduziblen Polynoms $m(X)=X^4-2$.
kann man einen Körper $\mathbb{Q}(\sqrt[4]{2})$ konstruieren, welcher
das Element enthält.

Auch der Körper $\mathbb{Q}(\!\sqrt{2})$ enthält die vierte Wurzel
$\sqrt[4]{2}$ nicht.
Um einen Körper zu konstruieren, der auch $\sqrt[4]{2}$ enthält, muss
$\mathbb{Q}(\!\sqrt{2})$ um die vierte Wurzel erweitert werden.
Dazu wird das Minimalpolynom von $\sqrt[4]{2}$ benötigt, welches
irreduzibel sein muss.
Das bei der Konstruktion von $\mathbb{Q}(\sqrt[4]{2})$ verwendete
Polynom $X^4-2$ ist jedoch nicht geeignet, da es über dem aktuelle
Koeffizientenkörper $\mathbb{Q}(\!\sqrt{2})$ nicht irreduzibel ist,
sondern die Faktorzerlegung
\[
X^4 - 2 = (X + \!\sqrt{2})(X - \!\sqrt{2})
\]
hat.
Beide Faktoren haben Koeffizienten in $\mathbb{Q}(\!\sqrt{2})$,
aber nicht alle Koeffizienten sind auch in $\mathbb{Q}$.
Einen Körper, der auch $\sqrt[4]{2}$ enthält, kann man also durch
Erweiterung von $\mathbb{Q}(\!\sqrt{2})$ mithilfe des Polynoms
$X^2-\!\sqrt{2}$ erhalten.
Man könnte diesen Körper als $\mathbb{Q}(\!\sqrt{2})(\sqrt[4]{2})$
schreiben.

Umgekehrt ist die Quadratwurzel $\!\sqrt{2}$ bereits in
$\mathbb{Q}(\sqrt[4]{2})$, da $\!\sqrt{2}=(\sqrt[4]{2})^2$.
Es ist also bereits $\!\sqrt{2}\in\mathbb{Q}(\sqrt[4]{2})$.
Es folgt, dass $\mathbb{Q}(\!\sqrt{2})(\sqrt[4]{2})
=
\mathbb{Q}(\sqrt[4]{2})$.
Die Erweiterung von $\mathbb{Q}$ erst um $\!\sqrt{2}$ 
und dann um $\sqrt[4]{2}$ liefert zwei verschiedene Körper,
während die Erweiterung um $\sqrt[4]{2}$ sofort den grossen
Körper liefert.

Anders sieht es bei der Erweiterung von $\mathbb{Q}$ um $\!\sqrt{2}$
und $\!\sqrt{3}$ aus.
Als Minimalpolynom von $\!\sqrt{3}$ kann $X^2-3$ verwendet werden,
welches auch irreduzibel ist über $\mathbb{Q}(\!\sqrt{2})$.
Die Erweiterung von $\mathbb{Q}(\!\sqrt{2})$ um $\!\sqrt{3}$
liefert den Körper
\begin{align*}
\mathbb{Q}(\!\sqrt{2})(\!\sqrt{3})
&=
\Bigl\{
a+b\!\sqrt{3}
\;
\Big|
\;
a,b\in\mathbb{Q}(\!\sqrt{2})
\Bigr\}
\\
&=
\Bigl\{
a_0+a_1\!\sqrt{2}+a_2\!\sqrt{3}+a_3\!\sqrt{6}
\;
\Big|
\;
a_k\in\mathbb{Q}
\Bigr\}.
\end{align*}
Die umgekehrte Reihenfolge führt auf
\begin{align*}
\mathbb{Q}(\!\sqrt{3})(\!\sqrt{2})
&=
\Bigl\{
a+b\!\sqrt{2}
\;
\Big|
\;
a,b\in\mathbb{Q}(\!\sqrt{3})
\Bigr\}
\\
&=
\Bigl\{
b_0+b_1\!\sqrt{3}+b_2\!\sqrt{2}+b_3\!\sqrt{6}
\;
\Big|
\;
b_k\in\mathbb{Q}
\Bigr\}.
\end{align*}
Offensichtlich entsteht der gleiche Körper.
Es kommt also nicht auf die Reihenfolge an und es ist zulässig,
\[
\mathbb{Q}(\!\sqrt{2})(\!\sqrt{3})
=
\mathbb{Q}(\!\sqrt{2},\!\sqrt{3})
=
\mathbb{Q}(\!\sqrt{3},\!\sqrt{2})
=
\mathbb{Q}(\!\sqrt{3})(\!\sqrt{2})
\]
zu schreiben.

%
% Transzendente Erweiterungen
%
\subsubsection{Transzendente Erweiterungen}
Da die Menge aller rationalen Polynome abzählbar unendlich ist,
kann es nur abzählbar viele über $\mathbb{Q}$ algebraische 
Zahlen geben.
Die Menge $\mathbb{R}$ der reellen Zahlen ist aber überabzählbar
unendlich, so dass die meisten reellen Zahlen nicht algebraisch
sind.
Sie heissen transzendent.

\begin{definition}[transzendent]
Ein Element $\alpha$ heisst \emph{transzendent} über $\mathbb{Q}$,
wenn es nicht algebraisch ist.
\index{transzendent}%
\end{definition}

Die Zahlen $\pi$ und $e$ sind beide transzendent.
Der Nachweis, dass eine Zahl transzendent ist, ist oft schwierig.
Schon 1844 hat Jospeh Liouville vermutet, dass $e$ transzendent 
\index{Liouville, Joseph}%
ist, der Nachweis ist aber Charles Hermite erst 1873 gelungen
\cite{buch:hermitebeweis}.
\index{Hermite, Charles}%
Für $\pi$ ist der Beweis deutlich anspruchsvoller.
Nach einem Satz von Lindemann und Weierstrass
\index{Lindemann, Charles}%
\index{Weierstrass, Karl}%
\cite{buch:lindemann-weierstrass} sind für verschiedene
algebraische Zahlen $\alpha_1,\dots,\alpha_n$
die Zahlen $e^{\alpha_1},\dots,e^{\alpha_n}$ linear unabhängig
über $\mathbb{Q}$.
Die Zahl $0$ ist sicher algebraisch.
Wäre auch $\pi$ eine algebraische Zahl, dann wäre auch $i\pi$ algebraisch.
Nach dem Satz von Lindemann-Weierstrass müssten dann $e^0=1$ und
$e^{i\pi}$ über den algebraischen Zahlen linear unabhängig sein,
was sie wegen der eulerschen Identität
\index{eulersche Identität}%
\index{Identität, eulersch}%
\[
e^{i\pi}+1=0
\]
offensichtlich nicht sind.
Der Widerspruch zeigt, dass $\pi$ transzendent sein muss.

%
% Algebraische Erweiterungen von F_p
%
\subsubsection{Algebraische Erweiterungen von $\mathbb{F}_p$}
Die Konstruktion der reellen Zahlen mithilfe von Cauchy-Folgen ist
nur möglich, weil es auf der Menge $\mathbb{Q}$ der rationalen
Zahlen eine Ordnungsrelation gibt, die der Idee des Abstands
$|a_n-a_m|$ von Folgengliedern einen Sinn gibt.
Die Definition eines algebraischen Elementes hingegen ist auch dann 
möglich, wenn es keine solche Ordnungsrelation gibt.
Die endlichen Körper $\mathbb{F}_p$ für Primzahlen $p$ sind von
dieser Art.

\begin{definition}[endliche Körper $\mathbb{F}_p$]
\index{endliche Körper}%
\index{Korper@Körper!endlich}%
Sei $p$ eine Primzahl.
Die Menge
\[
\mathbb{F}_p
=
\{ 0, \dots , p-1 \}
\]
der ganzzahligen Reste bei Division durch $p$
heisst der \emph{endlich Körper} der Ordnung $p$.
\index{endlicher Körper}%
\index{Körper, endlich}%
\index{Fp@$\mathbb{F}_p$}%
\end{definition}

In $\mathbb{F}_p$ sind wie in $\mathbb{Q}$ alle vier Grundoperationen
wohldefiniert, insbesondere auch das reziproke Element.
Diese letzte Aussage ist gleichbedeutend damit, dass die Multiplikation
\[
x\mapsto a\cdot x
\]
mit einem Element $a\in\mathbb{F}_p\setminus\{0\}$ eine Bijektion ist.
Da dies eine lineare Abbildung ist, muss man nur prüfen, ob es ein
Element $x\ne 0$ gibt derart, dass $ax=0$ in $\mathbb{F}_p$ ist.
Dies ist gleichbedeutend damit, dass 
\[
a\cdot n \equiv 0 \mod p
\qquad\Leftrightarrow\qquad
p \mid an
\]
gilt.
Da weder $a$ noch $n$ von $p$ geteilt werden können und $p$ eine
Primzahl ist, kann es eine solche Zahl $n$ nicht geben.

Da alle ganzen Zahlen, die sich nur um ein Vielfaches von $p$
unterscheiden, das gleiche Element in $\mathbb{F}_p$ darstellen,
kann es keine mit den arithmetischen Operationen verträgelich
Ordnungsrelation geben.
Schon die Unterschiedung zwischen positiven und negativen
Zahlen ist sinnlos, da jedes Element $a$ mit $0\le a<p$ auch
die äquivalente Darstellung $a-p$ hat, welche negativ ist.
Es folgt, dass es auch nicht sinnvoll sein kann, in $\mathbb{F}_p$
von Folgen und Grenzwerten zu sprechen.
So etwas wie eine Vervollständigung von $\mathbb{F}_p$ wie
die Vervollständigung von $\mathbb{Q}$ zu $\mathbb{R}$ kann es
also nicht geben.

Trotzdem ist es sinnvoll, von algebraischen Erweiturungen zu sprechen.
Wir betrachten dazu statt Polynomen über $\mathbb{Q}$ solche über
$\mathbb{F}_p$.
Wieder setzen wir
\[
\mathbb{F}_p[X]
=
\{
X^n + a_{n-1}X^{n-1} + \dots + a_1X + a_0
\mid
a_0,\dots,a_{n-1}\in \mathbb{F}_p
\}.
\]
Wie im Falle von rationalen Koeffizienten ist ein irreduzibles
Polynom eines, welches sich nicht als Produkt zweier Polynome von
positivem Grad schreiben lässt.
Dabei kann es aber durchaus passieren, dass ein Polynom, welches
über $\mathbb{Q}$ irreduzibel war, als Polynom über $\mathbb{F}_p$
betrachtet nicht mehr irreduzibel sein muss.

\begin{beispiel}
Für $p=5$ ist
\[
(X+3)(X+2)
=
X^2 + (3+2)X + 6
=
X^2 + 1
\in
\mathbb{F}_5[X].
\]
Das Polynom $X^2+1\in\mathbb{F}_5[X]$ ist also nicht irreduzibel,
im Gegensatz zum Polynom $X^2+1\in\mathbb{Q}[X]$, welches tatsächlich
irreduzibel ist.
Dies zeigt auch, dass im Gegensatz zu $\mathbb{Q}$ oder $\mathbb{R}$,
die Gleichung $X^2+1=0$ in $\mathbb{F}_p$ eine Lösung hat.
\end{beispiel}

Wie für rationale Koeffizienten lassen sich auch für ein algebraisches
Element über $\mathbb{F}_p$ die arithmetischen Grundoperationen
definieren.

\begin{beispiel}
\label{buch:komplex:zahlen:bsp:F5}
Für $p=5$ ist $X^2-2=X^2+3\in\mathbb{F}_5[X]$ ein irreduzibles Polynom,
wie man durch Durchprobieren der möglichen Faktoren $(X-a)(X-b)$ mit
$a,b\in\mathbb{F}_5$ nachprüfen kann.
Sei jetzt $\alpha$ ein Objekt, welches $\alpha^2=2$ erfüllt.
Die Rechenoperationen sind ganz analog wie im Falle rationaler
Koeffizienten definiert.
\end{beispiel}

Die Nullstelle $\alpha$ des Polynoms $X^2-2\in\mathbb{F}_4[X]$,
die in Beispiel~\ref{buch:komplex:zahlen:bsp:F5} zu $\mathbb{F}_5$
hinzugefügt worden ist, ist eine Quadratwurzel von $2$ in $\mathbb{F}_5$,
hat aber nichts mit der vertrauten Quadratwurzeln
$\!\sqrt{5}\approx 1.414213562373$ in $\mathbb{R}$ zu tun.
Während es für reelle Nullstellen von Polynomen in $\mathbb{Q}[X]$ 
immer eine Dezimaldarstellung gibt, gibt es keine alternative Darstellung
für $\alpha$.

\begin{beispiel}
Das Polynom $m(X)=X^3-2=X^3+5\in\mathbb{F}_7[X]$ ist irreduzibel, denn
wenn es faktorisiert werden könnte, müsste notwendigerweise einer der
Faktoren den Grad 1 haben.
Ein solcher Faktor wäre also von der Form $X-a$, wobei $a\in\mathbb{F}_7$
sein müsste.
Das Polynom müsste also eine Nullstelle in $\mathbb{F}_7[X]$ haben.
Die dritten Potenzen aller Elemente von $\mathbb{F}_7$ sind allerdings
\[
\begin{tabular}{|>{$}l<{$}|>{$}c<{$}>{$}c<{$}>{$}c<{$}>{$}c<{$}>{$}c<{$}>{$}c<{$}>{$}c<{$}|}
\hline
a   & 0 & 1 & 2 & 3 & 4 & 5 & 6 \\
\hline
a^3 & 0 & 1 & 1 & 6 & 1 & 1 & 6 \\
\hline
\end{tabular}
\]
Keine dritte Potenz ergibt 2.

Fügt man $\mathbb{F}_7$ das Element $\alpha$, welches eine Nullstelle
von $m$ ist, zu $\mathbb{F}_7$ hinzu, erhält man neu die Menge
$\mathbb{F}_7(\alpha)$.
Lässt man Koeffizienten aus $\mathbb{F}_7(\alpha)$ zu, lässt sich das
Polynom $m$ wie in Beispiel~\ref{buch:komplex:zahlen:bsp:kubikwurzel}
als
\[
x^3-2
=
(x-\alpha)(\underbrace{x^2+\alpha x + \alpha^2}_{\displaystyle = b})
\]
faktorisieren.
Auch in diesem Beispiel ist eine weitere Erweiterung
$\mathbb{F}_7(\alpha,\beta)$ von $\mathbb{F}_p(\alpha)$ durch eine
Nullstelle $\beta$ des Polynoms $b\in\mathbb{F}_7(\alpha)$ nötig,
damit die anderen zwei Nullstellen von $m$ ebenfalls gefunden werden
können.
\end{beispiel}

Die Adjunktion einer Nullstelle eines irreduziblen Polynoms 
\index{Adjunktion}%
aus $\mathbb{F}_p[x]$ zu $\mathbb{F}_p$ erzeugt also wieder
die Menge 
\[
\mathbb{F}_p(\alpha)
=
\{
a_0+a_1\alpha + \dots a_{n-1}\alpha^{n-1}
\mid
a_0,\dots,a_{n-1}\in\mathbb{F}_p
\},
\]
die ein Vektorraum über $\mathbb{F}_p$ ist, mit den $n$ Monomen
\[
1, \alpha, \dots, \alpha^{n-1}
\]
als Basis.
$\mathbb{F}_p(\alpha)$ enthält somit $p^n$ Elemente.

\begin{definition}[endliche Körper]
Die Menge $\mathbb{F}_p(\alpha)$ mit $p^n$ Elementen heisst der
\emph{endliche Körper mit $p^n$} Elementen.
\index{endliche Körper}%
\index{Korper@Körper!endlich}%
\end{definition}

Fügt man einem Körper mit $p^n$ Elementen ein weiteres Element mit
einem Minimalpolynom vom Grad $l$ hinzu, entsteht eine Menge mit
mit $(p^n)^l=p^{nl}$ Elementen.
Man kann zeigen, dass alle Körper mit $p^n$ wie man sagt isomorph sind,
ganz unabhängig davon, was für ein Minimalpolynom verwendet wird
und schreibt daher unabhängig vom Minimalpolynom
$\mathbb{F}_p(\alpha)=\mathbb{F}_{p^n}$.
Die endlichen Körper spielen eine bedeutende Rolle in der Kryptographie.
Zum Beispiel verwendet die Spezifikation des AES-Algorithmus die 
\index{AES-Algorithmus}%
Arithmetik in einem Körper $\mathbb{F}_{2^8}$ mit $2^8$ Elementen,
die im Code als Bytes dargestellt werden.

%
% Die ganzen gaussschen Zahlen Z[i]
%
\subsubsection{Die ganzen gaussschen Zahlen $\mathbb{Z}[i]$}
\input{chapters/010-komplex/fig/fig-gausszahlen.tex}%
Die Erweiterung einer Zahlenmenge um ein neues Element, welches als
Nullstelle eines irreduziblen Polynoms definiert wurde, lässt sich
auch auf die Menge $\mathbb{Z}$ anwenden.
Für das Polynom $m(X) = X^2+1\in\mathbb{Z}[X]$ wird die Nullstelle
mit $i$ bezeichnet.
Die Rechenregeln folgen unmittelbar aus den Rechenregeln für
Polynome mit ganzzahligen Koeffizienten durch Anwendung der
Identität $i^2+1=0$ oder $i^2=-1$.
Es entsteht die Menge der ganzen gaussschen Zahlen gemäss der
folgenden Definition.

\begin{definition}[ganze gausssche Zahlen]
\index{ganze gausssche Zahlen}%
Die Elemente
\[
\mathbb{Z}[i]
=
\{
a+bi
\mid
a,b\in\mathbb{Z}
\}
\]
heisst die Menge der ganzen gaussschen Zahlen.
Sie bildet einen kommutativen Ring mit Eins, der den Ring $\mathbb{Z}$
der ganzen Zahlen als echten Unterring enthält.
Die Komponente $a$ heisst der \emph{Realteil} von $z=a+bi$, geschrieben
$a=\operatorname{Re}z$, $b=\operatorname{Im}z$ heisst Imaginärteil.
\end{definition}

Die ganzen gaussschen Zahlen können dargestellt werden als die
Gitterpunkt in einer Ebene, deren horizontale Achse der reelle
Zahlenstrahl ist, während die vertikale Achse die imaginären
Zahlen darstellt (Abbildung~\ref{buch:komplex:zahlen:fig:gausszahlen}).

Die Ringoperationen sind
\begin{align*}
\text{Addition:}&&
(a+bi)\pm(c+di)
&=
(a\pm c) + (b\pm d)i
\\
\text{Multiplikation:}&&
(a+bi)(c+di)
&=
ac-bd
+
(ad+bc)i.
\end{align*}
Nicht jede ganze gausssche Zahl lässt sich invertieren.
Die Inverse wird durch Erweiterung mit $a-bi$ als
\[
\frac{1}{a+bi}\cdot\frac{a-bi}{a-bi}
=
\frac{a-bi}{a^2+b^2}
=
\frac{a}{a^2+b^2}
-
\frac{b}{a^2+b^2}i
\]
gefunden.
Die Inverse existert also genau dann, wenn $a^2+b^2=1$ ist.
Dies ist nur für die Zahlen $\pm 1$ und $\pm i$ der Falle.

\begin{definition}[Betrag]
Für eine ganze gausssche Zahl $z=a+bi\in\mathbb{Z}[i]$ heisst die
Zahl $|z|$ definiert durch $|z|^2 = a^2+b^2$ der \emph{Betrag} von $z$.
\index{Betrag}%
\end{definition}

Für zwei das Produkt $z=(a+bi)(c+di)$ zweier ganzer gaussscher Zahlen gilt
\begin{align*}
|z|^2
&=
(ac-bd)^2 + (ad+bc)^2
\\
&=
a^2c^2-2abcd+b^2d^2
+
a^2d^2
+
2abcd
+
b^2c^2
\\
&=
a^2c^2 + a^2d^2 + b^2c^2 + b^2d^2
\\
&=
(a^2+b^2)(c^2+d^2)
\\
&=
|a+bi|^2 |c+di|^2
\intertext{oder}
|z|
&=
|a+bi|\cdot |c+di|.
\end{align*}
Mit anderen Worten: Wenn sich eine ganze gausssche Zahl faktorisieren
lässt, dann lassen sich auch die Beträge faktorisieren.

Tatsächlich kann man im Ring $\mathbb{Z}[i]$ der ganzen gausschen Zahlen
eine Theorie der Faktorisierung aufstellen, die ganz analog ist zur
Theorie der Primfaktoren in $\mathbb{Z}$.


\input{chapters/010-komplex/13-komplex.tex}
%
% 14-matrix.tex -- Matrixschreibweise für komplexe Zahlen
%
% (c) 2026 Prof Dr Andreas Müller
%
\subsection{Matrixschreibweise für komplexe Zahlen
\label{buch:komplex:subsection:matrixschreibweise}}
In der Konstruktion der Körpererweiterung $\mathbb{Q}(\alpha)$
mit einer Nullstelle $\alpha$ des Minimalpolynoms $m(X)\in\mathbb{Q}[X]$ 
wurde nie die Frage gestellt, woher $\alpha$ eigentlich kommt.
Für $m=X^2-2$ kann sich $\alpha$ als die Quadratwurzel
$\!\sqrt{2}\in\mathbb{R}$ vorstellen, aber für die Körpererweiterung
von $\mathbb{F}_5$ um eine Nullstelle von $X^2-2\in\mathbb{F}_5[X]$
gibt es keinen solchen Überkörper, der alle denkbaren Wurzeln
enthält.
Und bei der Erweiterung $\mathbb{R}\subset \mathbb{R}(i)=\mathbb{C}$
wird $\mathbb{R}$ die imaginäre Einheit $i$ hinzugefügt, für die es
ebenfalls keine Realisierung in $\mathbb{R}$ gibt.
Schlimmer noch: die imaginäre Einheit hat Eigenschaften, die den 
Eigenschaften aller bisher bekannten Zahlen widerspricht, nämlich
die Eigenschaft, dass das Quadrat einer Zahl immer $\ge 0$ ist.

So schreibt auch Leonhard Euler in seiner \emph{Vollständigen
Anleitung zur Algebra}
\index{Euler, Leonhard}%
\cite[S.~61,  Abschnitt~144--145]{buch:euler-algebra}:
\begin{quote}
Daher bedeuten alle diese Ausdrücke $\!\sqrt{-1}$, $\!\sqrt{-2}$,
$\!\sqrt{-3}$, $\!\sqrt{-4}$, etc. solche unmögliche oder
Imaginärzahlen, weil dadurch Quadratwurzeln von Negativzahlen
angezeigt werden.

Von diesen behauptet man also mit allem Rechte, dass sie weder
grösser oder kleiner sind als nichts; und auch nicht einmal nichts
selbst, als aus welchem Grund sie folglich für unmöglich gehalten
werden müssen.

Gleichwohl aber werden sie unserm Verstande dargestellet, und finden
in unserer Einbildung statt; daher sie auch bloss eingebildete
Zahlen genennet werden.
Ungeachtet aber diese Zahle, als z.~B. $\!\sqrt{-4}$, ihrer Natur
nach ganz und gar unmöglich sind, so haben wir davon doch einen
hinlänglichen Begriff, indem wir wissen, dass dadurch eine solche
Zahl angedeutet werde, welche mit sich selbst multiplicirt zum
Producte -4 hervorbringe; und dieser Begriff ist zureichend um diese
Zahlen in der Rechnung gehörig zu behandeln.
\end{quote}
Euler war sich also des Widerspruchs durchaus bewusst und versucht
auch gar nicht, eine ``physikalische'' Rechtfertigung für die
imaginären Zahlen zu konstruieren.
Vielmehr weisst er darauf hin, dass es genügt, dass wir mit diesen
Zahlen ohne Widerspruch algebraisch arbeiten können, so wie wir
das bei der Konstruktion der algebraischen Körpererweiterung getan
haben.

Trotzdem muss man die Frage stellen, wie Euler sicher sein konnte,
dass die Algebra der komplexen Zahlen wirklich widerspruchsfrei ist.
Seine umfassende Erfahrung und Intuition in der Algebra dieser Zahlen
wird ihn wohl zu dieser Annahme verleitet haben, ohne dass er einen
weiteren Gedanken darauf verwendet, wie man diese Widerspruchsfreiheit
beweisen könnte.
Zugegebenermassen ist dies für die komplexen Zahlen auch nicht eine
grosse Schwierigkeit, da sich mit der Darstellung von $\mathbb{C}$
als komplexe Zahlenebene eine Konstruktion als Paare von reellen
Zahlen mit Rechenoperationen, für die man die Axiome eines Körpers
verifizieren kann, aufdrängt.

Etwas weniger klar ist die Situation jedoch bei algebraischen
Erweiterungen von $\mathbb{Q}$ und endlichen Körpern.
Zwar ist für Nullstellen von quadratischen Polynomen
immer noch die Vorstellung einer Ebene angemessen, für Körpererweiterungen
um Nullstellen von Polynomen beliebigen Grades funktioniert diese
Intuition jedoch nicht mehr.
Gesucht ist daher eine allgemeine Konstruktion eines Modells für 
algebraische Körpererweiterungen $K(\alpha)$ für Nullstellen beliebiger
Minimalpolynome $m(X)\in K[X]$.
Im Folgenden soll gezeigt werden, wie ein solches Modell innerhalb
der Matrizenalgebra konstruiert werden kann.

\subsubsection{Algebraische Erweiterungen als Matrizen}
Wir beginnen mit einem Beispiel.
Sei $A$ die Matrix
\begin{equation}
A 
=
\begin{pmatrix}
0&2\\
1&0
\end{pmatrix}
\in
M_{2\times 2}(\mathbb{Q}).
\label{buch:komplex:zahlen:eqn:Awurzel2}
\end{equation}
Das Quadrat dieser Matrix ist
\[
A^2
=
\begin{pmatrix}
0&2\\
1&0
\end{pmatrix}
\begin{pmatrix}
0&2\\
1&0
\end{pmatrix}
=
\begin{pmatrix}
2&0\\
0&1
\end{pmatrix}
=
2I
\qquad\text{oder}\qquad
A^2-2I = 0,
\]
wobei $I$ die Einheitsmatrix ist.
Mit $A$ ist also ein algebraisches Objekt gefunden, welches die
Polynomgleichung $x^2-2=0$ erfüllt, sofern man den konstanten
Term des Polynoms auf der linken Seite als $2A^0=2I$ liest und 
die 0 auf der rechten Seite als die Nullmatrix.

Die inverse Matrix von $A$ ist
\[
A^{-1}
=
\begin{pmatrix}
0       & 1 \\
\frac12 & 0
\end{pmatrix}
=
\frac12 A,
\]
was sich mit der früher gefundenen Formel für die Inverse eines
algebraischen Elementes deckt.
Die negative Matrix $-A$ ist eine weitere Nullstelle der Gleichung
$x^2-2=0$.

Die Konstruktion ist natürlich auf den Körper $\mathbb{Q}$ oder
die Konstante $2$ beschränkt.
Die Matrix 
\[
A
=
\begin{pmatrix}
0 & a \\
1 & 0
\end{pmatrix}
\qquad\text{mit Inverser}\qquad
A^{-1}
=
\begin{pmatrix}
0      & 1 \\
a^{-1} & 0
\end{pmatrix}
=
a^{-1}A
\]
ist eine Nullstelle der Gleichung $x^2-a=0$ für ein beliebiges 
invertierbares Element $a\in K^*$ des Koeffizientenkörpers.
$-A$ ist eine weitere solche Lösung.

Die Beispiele zeigen, dass die Matrizenalgebra offenbar flexibel
genug ist, um für beliebige Quadratwurzeln eine Matrixdarstellung
zu finden.
Die Darstellung ist allerdings nicht eindeutig, wie der Versuch
der Lösung der quadratische Gleichung
\[
x^2 + px + q = 0
\]
zeigt.
Nach der bekannten Lösungsformel erwarten wir die Lösung in der
Form
\[
x_{\pm}
=
-\frac{p}2 \pm \!\sqrt{\rlap{$\phantom{\bigg|}$}\smash{\biggl(\frac{p}{2}\biggr)^2 - q}}.
\]
Übersetzt in die Matrixform müsste 
\[
X_{\pm}
=
-\frac{p}2I
\pm
\begin{pmatrix}
0 & \sqrt{(p/2)^2-q} \\
1 & 0
\end{pmatrix}
\]
eine Lösung der Matrixgleichung 
\begin{equation}
X^2 + pX + q = 0
\label{buch:komplex:zahlen:eqn:pqgleichung}
\end{equation}
sein.
Wir berechnen dazu das Produkt
\begin{align*}
X_{\pm}^2
&=
\begin{pmatrix}
(p/2)^2 - q & 0 \\
      0     & (p/2)^2 - q
\end{pmatrix}
=
((p/2)^2-q) I
\end{align*}
und setzen in die Gleichung
\[
X_\pm^2 + pX + q
=
((p/2)^2-q)I
+
pX_\pm
+
qI
=
\begin{pmatrix}
p^2/4 & q + p^3/4 -pq \\
p + q  & p^2/4
\end{pmatrix}
\]
ein.
Die Gleichung ist also nicht erfüllt.
Der Grund ist, dass es in der Menge $M_{2\times 2}(K)$ unendlich
viele Quadratwurzeln gibt.
Ist $T\in \operatorname{GL}_2(K)$ eine beliebige invertierbare 
$2\times 2$-Matrix, dann gilt für $A'=TAT^{-1}$
\[
A^{\prime 2}
=
(TAT^{-1})(TAT^{-1})
=
TA(\underbrace{T^{-1}T}_{\displaystyle=I})AT^{-1}
=
TA^2T^{-1}.
\]
Ist $A$ eine Quadratwurzel von $a$, dann ist
\[
A^{\prime 2}
=
TA^2T^{-1}
=
T(aI)T^{-1}
=
aI,
\]
$A'$ ist also auch eine Quadratwurzel von $a$.
Die Matrizenalgebra ist also viel grösser als zur Darstellung von
Quadratwurzeln nötig ist, zu gross um eine Quadratwurzel eindeutig 
festzulegen.

%
% Eine Matrix als Nullstelle des Minimalpolynoms
%
\subsubsection{Eine Matrix als Nullstelle des Minimalpolynoms}
Es lässt sich aber trotzdem eine Matrix angeben, die als Lösung
der Gleichung \eqref{buch:komplex:zahlen:eqn:pqgleichung} angesehen
werden kann, nämlich
\[
X
=
\begin{pmatrix}
-p & -q \\
 1 &  0
\end{pmatrix}.
\]
Tatsächlich ist
\begin{align*}
X^2 + pX + qI
&=
\begin{pmatrix}
-p & -q \\
 1 &  0
\end{pmatrix}^2
+
p
\begin{pmatrix}
-p & -q \\
 1 &  0
\end{pmatrix}^2
+qI
\\
&=
\begin{pmatrix}
p^2 - q & pq \\
   -p   & -q
\end{pmatrix}
+
\begin{pmatrix}
-p^2 & -pq \\
 p   &  0 
\end{pmatrix}
+
\begin{pmatrix}
q & 0 \\
0 & q
\end{pmatrix}
=
\begin{pmatrix}
0 & 0 \\
0 & 0
\end{pmatrix}.
\end{align*}
Die Matrix $X$ ist also eine Lösung der quadratischen Gleichung
\eqref{buch:komplex:zahlen:eqn:pqgleichung}.

Die Konstruktion des letzten Beispiels lässt sich verallgemeinern
zum folgenden Satz, der besagt, dass jedes über $K$ algebraische
Element vom Grad $n$ eine Matrixdarstellung in $M_{n\times n}(K)$
hat, die Nullstelle der Polynomgleichung ist, die das algebraische
Element definiert.

\begin{satz}
\label{buch:komplex:zahlen:satz:matrixalgebraisch}
Ist $a(X) = X^n + a_{n-1}X^{n-1}+\dots + a_1X+a_0 \in K[X]$ ein
normiertes Polynom vom Grad $n$ mit Koeffizienten im Körper $K$,
dann ist die Matrix
\[
A
=
\begin{pmatrix}
-a_{n-1} & -a_{n-2} & \dots  &  -a_1   &  -a_0   \\
    1    &     0    & \dots  &    0    &    0    \\
    0    &     1    & \dots  &    0    &    0    \\
 \vdots  &  \vdots  & \ddots &  \vdots &  \vdots \\
    0    &     0    & \dots  &    1    &    0 
\end{pmatrix}
\]
eine Nullstelle des Polynoms, also $a(A)=0$.
\end{satz}

\begin{proof}
Wir berechnen als erstes das charakteristische Polynom von $X$ mithilfe
des laplaceschen Entwicklungssatzes, indem wir nach der ersten Spalte
entwickeln.
Zur Vereinfachung schreiben wir
\[
A_{k}
=
\begin{pmatrix}
    -a_{k-1}     & -a_{k-2} & \dots  &  -a_1   &  -a_0   \mathstrut\\
        1        & -\lambda & \dots  &    0    &    0    \mathstrut\\
        0        &     1    & \dots  &    0    &    0    \mathstrut\\
     \vdots      &  \vdots  & \ddots &  \vdots &  \vdots \\
        0        &     0    & \dots  &    1    & -\lambda\mathstrut
\end{pmatrix}
\in
M_{k\times k}(K).
\]
Für die Determinante von $A_k$ folgt durch Entwicklung nach der ersten
Spalte die Rekursionsformel
\begin{align}
\det A_k
&=
-a_{k-1}
\left|
\begin{matrix}
 -\lambda &     0    & \dots  &    0     &    0    \mathstrut\\
     1    & -\lambda & \dots  &    0     &    0    \mathstrut\\
  \vdots  &  \vdots  & \ddots &  \vdots  &  \vdots \\
     0    &     0    & \dots  & -\lambda &    0    \mathstrut\\
     0    &     0    & \dots  &    1     & -\lambda\mathstrut
\end{matrix}
\right|
-
\left|
\begin{matrix}
 -a_{k-2} & -a_{k-3} & \dots  &  -a_1    &  -a_0   \mathstrut\\
     1    & -\lambda & \dots  &    0     &    0    \mathstrut\\
  \vdots  & \vdots   & \ddots &  \vdots  &  \vdots \\
     0    &    0     & \dots  & -\lambda &    0    \mathstrut\\
     0    &    0     & \dots  &    1     & -\lambda\mathstrut
\end{matrix}
\right|
\notag
\\
&=
-a_{k-1}(-\lambda)^{k-1}
-\det A_{k-2}
\label{buch:komplex:zahlen:eqn:xkrekursion}
\intertext{und die Anfangsbedingung}
\det A_2
&=
\biggl|
\begin{matrix}
-a_1 & -a_0     \mathstrut\\
  1  & -\lambda \mathstrut
\end{matrix}
\biggr|
=
a_1\lambda + a_0.
\label{buch:komplex:zahlen:eqn:xkstart}
\end{align}
Durch wiederholte Anwendung der Rekursionsformel
\eqref{buch:komplex:zahlen:eqn:xkrekursion}
ausgehend vom Anfangswert
\eqref{buch:komplex:zahlen:eqn:xkstart}
\begin{align*}
\det A_3
&=
-a_2(-\lambda)^2 - \det A_2
&&=
-a_2\lambda^2 -a_1\lambda - a_0
\\
\det A_4
&=
-a_3(-\lambda)^3 - \det A_3
&&=
a_3 \lambda^3 + a_2\lambda^2 + a_1\lambda + a_0
\\
\det A_5
&=
-a_4(-\lambda)^4 - \det A_4
&&=
-a_4\lambda^4
-a_3\lambda^3 - a_2\lambda^2 -a_1\lambda - a_0
\\
&\phantom{n}\vdots
\\
\det A_k
&=
-a_{k-1}(-\lambda)^{k-1}
-\det A_{k-1}
&&=
(-1)^k\bigl(
a_{k-1}\lambda^{k-1} + \dots + a_1\lambda + a_0
\bigr)
\end{align*}
findet man Ausdrücke für die Determinante $\det A_k$.

Damit kann man jetzt auch das charakteristische Polynom bestimmen.
Dieses ist
\begin{align}
\chi_A(\lambda)
&=
\det (A-\lambda I)
\notag
\\
&=
(-a_{n-1}-\lambda)
\left|
\begin{matrix}
 -\lambda & \dots  &    0    &    0    \mathstrut\\
     1    & \dots  &    0    &    0    \mathstrut\\
  \vdots  & \ddots &  \vdots &  \vdots \mathstrut\\
     0    & \dots  &    1    & -\lambda\mathstrut
\end{matrix}
\right|
-
\left|
\begin{matrix}
-a_{n-2} & -a_{n-3} & \dots  &  -a_1   &  -a_0   \mathstrut\\
    1    & -\lambda & \dots  &    0    &    0    \mathstrut\\
 \vdots  &  \vdots  & \ddots &  \vdots &  \vdots \mathstrut\\
    0    &     0    & \dots  &    1    & -\lambda\mathstrut
\end{matrix}
\right|
\notag
\\
&=
(-a_{n-1}-\lambda)(-\lambda)^{n-1}
-
\det A_{n-2}
\\
&=
(-1)^{n-2}
\bigl(
\lambda^n
+
a_{n-1}\lambda^{n-1}
\bigr)
-
(-1)^{n-1}
\bigl(
a_{n-2}\lambda^{n-2} + \dots + a_1\lambda + a_0
\bigr)
\notag
\\
&=
(-1)^n
\bigl(
\lambda^n
+
a_{n-1}\lambda^{n-1}
+
a_{n-2}\lambda^{n-2}
+
\dots
+
a_1\lambda
+
a_0
\bigr)
\notag
\\
&=
(-1)^n a(\lambda).
\label{buch:komplex:zahlen:eqn:charmin}
\end{align}
Das charakteristische Polynom ist also bis auf ein Vorzeichen das
ursprüngliche Polynom $a(X)\in K[x]$.

Der Satz von Cayley-Hamilton \cite[Satz 13.23]{buch:linalg} besagt, dass 
\index{Cayley-Hamilton, Satz von}%
\index{Satz!von Cayley-Hamilton}%
jede Matrix Nullstelle ihres charakteristischen Polynoms ist, dass also
$\chi_A(A)=0$ ist.
Es folgt daher aus \eqref{buch:komplex:zahlen:eqn:charmin},
dass auch $m(X)=0$ ist.
Damit ist der Satz vollständig bewiesen.
\end{proof}

\begin{beispiel}
Für das irreduzible Polynom $a(X) = X^2-2$ ist
$a_1=0$ und $a_0=-2$.
Die Formel des Satzes liefert daher für 
\[
A
=
\begin{pmatrix}
-a_1&-a_0\\
  1 &  0
\end{pmatrix}
=
\begin{pmatrix}
0 & 2 \\
1 & 0
\end{pmatrix},
\]
also genau die Matrix, die in
\eqref{buch:komplex:zahlen:eqn:Awurzel2}
als Quadratwurzel von $2$ vorgeschlagen wurde.
\end{beispiel}

%
% Die imaginäre Einheit als Matrix
%
\subsubsection{Die imaginäre Einheit als Matrix}
Die imaginäre Einheit ist Nullstelle des irreduziblen Polynoms
$X^2+1$.
Nach Satz~\ref{buch:komplex:zahlen:satz:matrixalgebraisch} kann
die imaginäre Einheit durch die $2\times 2$-Matrix
\[
J
=
\begin{pmatrix*}[r]
0 & -1 \\
1 &  0
\end{pmatrix*}
\]
dargestellt werden.
Tatsächlich sind die Potenzen von $J$
\begin{align*}
J^2
&=
\begin{pmatrix*}[r]
-1& 0\\
 0&-1
\end{pmatrix*}
=
-I
\\
J^3
&=
J^2J
=
-IJ
=
-J
\\
J^4
&=
J^2J^2
(-I)(-I)
=
I.
\\
J^k
&=
\begin{cases}
(-1)^{k/2}I&\qquad \text{$k$ gerade} \\
(-1)^{(k-1)/2}J&\qquad \text{$k$ ungerade,} \\
\end{cases}
\end{align*}
wie man dies von $i$ erwartet.

Die komplexe Zahl $z=a+bi$ wird also durch die Matrix
\[
Z = aI+bJ
=
\begin{pmatrix*}[r]
a&-b\\
b& a
\end{pmatrix*}
\]
dargestellt.
Wir rechnen nach, dass die inverse Matrix von $Z$ auch die
Matrix ist, die zur komplexen Zahl $z^{-1}$ gehört.
Die inverse Matrix ist
\begin{align*}
Z^{-1}
&=
\frac{1}{\det Z}
\begin{pmatrix*}[r]
 a & b \\
-b & a
\end{pmatrix*}
=
\frac{1}{a^2+b^2}
\begin{pmatrix*}[r]
 a & b \\
-b & a
\end{pmatrix*}
=
\frac{a}{a^2+b^2}I
+
\frac{-b}{a^2+b^2}J.
\end{align*}
Dies ist genau die Matrix, die zur reziproken komplexen Zahl
\[
\frac{1}{z}
=
\frac{a}{a^2+b^2}
+
\frac{-b}{a^2+b^2}i
\] 
gehört.

Man beachte, dass in dieser Matrixdarstellung der komplexen Zahlen
die erste Spalte ein Vektor ist, der die komplexe Zahl als Vektor
in der gaussschen Zahlenebene darstellt.

%
% 15-geometrisch.tex -- Geometrische Beschreibung der Rechenoperationen in C
%
% (c) 2026 Prof Dr Andreas Müller
%
\subsection{Geometrische Beschreibung der Rechenoperationen in $\mathbb{C}$}
\input{chapters/010-komplex/fig/fig-gaussebene.tex}%
Eine komplexe Zahl $z=a+bi$ kann sowohl als Vektor
\[
z = \begin{pmatrix}a\\b\end{pmatrix}
\]
in der gaussschen Zahlenebene der
\index{Zahlenebene}%
\index{Zahlenebene!gausssche}%
Abbildung~\ref{buch:komplex:zahlen:fig:gaussebene}
als auch als $2\times 2$-Matrix beschrieben werden.
Schreiben wir das Produkt $(a+bi)(c+di)$ in Matrix-Form als
\begin{align*}
(aI+bJ)(cI+dJ)
&=
\begin{pmatrix*}[r]
a&-b\\
b& a
\end{pmatrix*}
\begin{pmatrix*}[r]
c&-d\\
d& c
\end{pmatrix*}
\intertext{und betrachten wir nur die erste Spalte des Resultats,
finden wir das Produkt als Spaltenvektor}
\begin{pmatrix*}
ac-bd\\
bc+ad
\end{pmatrix*}
&=
\begin{pmatrix*}[r]
a&-b\\
b& a
\end{pmatrix*}
\begin{pmatrix*}[r]
c\\
d
\end{pmatrix*}
\end{align*}
Wie kann man dieses Produkt geometrisch in der gaussschen Zahlenebene
beschreiben?

Wir schreiben $r=\!\sqrt{a^2+b^2}$ und schreiben die Matrix von $a+bi$
als
\[
aI+bJ
=
r
\begin{pmatrix*}[r]
a/r & -b/r \\
b/r &  a/r
\end{pmatrix*}.
\]
Für die Matrixelemente gilt
\[
\biggl(\frac{a}{r}\biggr)^2
+
\biggl(\frac{b}{r}\biggr)^2
=
\frac{a^2}{r^2}
+
\frac{b^2}{r^2}
=
\frac{a^2+b^2}{r^2}
=
1.
\]
Es gibt daher einen Winkel $\varphi$ derart, dass
\[
\cos\varphi = \frac{a}{r}
\qquad\text{und}\qquad
\sin\varphi = \frac{b}{r}.
\]
Damit wird die Matrix zu
\[
aI+bJ
=
r
\begin{pmatrix*}[r]
\cos\varphi & -\sin\varphi \\
\sin\varphi &  \cos\varphi
\end{pmatrix*}
=
r\,D_\varphi,
\]
wobei $D_\varphi$ die Drehung der Ebene um den Winkel $\varphi$ im
Gegenuhrzeigersinn ist.
Die Multiplikation mit der komplexen Zahl $a+bi$ dreht also einen
Vektor in der gaussschen Zahlenebene um den Winkel $\varphi$ und
streckt ihn anschliessend um $r=|a+bi|$.

\begin{definition}[Polardarstellung komplexer Zahlen]
\index{Polardarstellung}%
Ist $z=a+bi$ eine komplexe Zahl, dann gibt es einen Winkel $\varphi$ derart,
dass 
\[
z
=
a + bi
=
r(\cos\varphi + i\sin\varphi)
\qquad\text{mit}\qquad
r=|z|.
\]
Der Winkel $\varphi$ heisst das \emph{Argument} $\varphi=\operatorname{arg}z$
der komplexen Zahl $z$.
\index{Argument einer komplexen Zahl}%
\index{argz@$\operatorname{arg}z$}%
\end{definition}


%
% 16-quaternionen.tex -- Quaternionen
%
% (c) 2026 Prof Dr Andreas Müller
%
\subsection{Quaternionen}
Die Matrix $J$ ist nicht die einzige Matrix mit der Eigenschaft
$J^2=-I$.
\index{Quaternionen}%
Lässt man komplexe Einträge in den Matrizen zu, haben auch die
Matrizen
\begin{align*}
-i\sigma_x
&=
-i
\begin{pmatrix*}[r]
0&1\\
1&0
\end{pmatrix*}
=
\begin{pmatrix*}[r]
 0&-i\\
-i& 0
\end{pmatrix*}
&&\Rightarrow&
(-i\sigma_x)^2
&=
\begin{pmatrix*}[r]
-1& 0\\
 0&-1
\end{pmatrix*}
=
-I
\\
-i\sigma_y
&=
-i
\begin{pmatrix*}[r]
0&-i\\
i& 0
\end{pmatrix*}
=
\begin{pmatrix*}[r]
0&-1\\
1&0
\end{pmatrix*}
=
J
&&\Rightarrow&
(-i\sigma_y)^2
&=
J^2
=
-I
\\
-i\sigma_z
&=
-i
\begin{pmatrix*}[r]
1&0\\
0&-1
\end{pmatrix*}
=
\begin{pmatrix*}[r]
-i&0\\
 0&i
\end{pmatrix*}
&&\Rightarrow&
(-i\sigma_z)^2
&=
\begin{pmatrix*}[r]
-1& 0\\
 0&-1
\end{pmatrix*}
=
-I.
\end{align*}
diese Eigenschaft.

\begin{definition}[Pauli-Matrizen]
Die Matrizen
\begin{align*}
\sigma_x
&=
\begin{pmatrix*}[r]
0&1\\
1&0
\end{pmatrix*},
&
\sigma_y
&=
\begin{pmatrix*}[r]
0&-i\\
i& 0
\end{pmatrix*}
&&\text{und}
&
\sigma_z
&=
\begin{pmatrix*}[r]
1&0\\
0&-1
\end{pmatrix*}
\end{align*}
heissen die \emph{Pauli-Matrizen}.
\index{Pauli-Matrix}%
\end{definition}

Die Pauli-Matrizen spielen eine wichtige Rolle in der Beschreibung des
Spins in der Quantenmechanik.
\index{Quantenmechanik}%
\index{Spin}%
Sie können auch verwendet werden, um die Gruppe $\operatorname{SU}(2)$ 
der spezielle unitären Matrizen zu beschreiben.
\index{unitär}%

Schreiben wir
\begin{align}
\boldsymbol{i}
&=
-i\sigma_x
=
\begin{pmatrix}
0&-i\\
-i&0
\end{pmatrix}
,
&
\boldsymbol{j}
&=
-i\sigma_y
=
\begin{pmatrix}
0&-1\\
1&0
\end{pmatrix}
=
J
&&\text{und}
&
\boldsymbol{k}
&=
-i\sigma_z
=
\begin{pmatrix}
-i & 0 \\
 0 & i
\end{pmatrix}
,
\label{buch:komplex:zahlen:eqn:quaternionen}
\end{align}
erhalten wir zusätzlich zu
$\boldsymbol{i}^2
=
\boldsymbol{j}^2
=
\boldsymbol{k}^2
=-1$
durch Nachrechnen die folgenden Beziehungen:
\begin{align}
\boldsymbol{i}\boldsymbol{j}
&=
\begin{pmatrix*}[r]
 0&-i\\
-i& 0
\end{pmatrix*}
\begin{pmatrix*}[r]
 0&-1\\
 1& 0
\end{pmatrix*}
=
\begin{pmatrix*}[r]
-i & 0 \\
 0 & i
\end{pmatrix*}
=
\boldsymbol{k},
&
\boldsymbol{j}\boldsymbol{i}
&=
\begin{pmatrix*}[r]
 0&-1\\
 1& 0
\end{pmatrix*}
\begin{pmatrix*}[r]
 0&-i\\
-i& 0
\end{pmatrix*}
=
\begin{pmatrix*}[r]
i & 0 \\
0 & -i
\end{pmatrix*}
=
-\boldsymbol{k},
\label{buch:komplex:zahlen:eqn:quatrel}
\\
\boldsymbol{j}\boldsymbol{k}
&=
\begin{pmatrix*}[r]
 0&-1\\
 1& 0
\end{pmatrix*}
\begin{pmatrix*}[r]
-i & 0 \\
 0 & i
\end{pmatrix*}
=
\begin{pmatrix*}[r]
 0&-i\\
-i& 0
\end{pmatrix*}
=
\boldsymbol{i},
&
\boldsymbol{k}\boldsymbol{j}
&=
\begin{pmatrix*}[r]
-i & 0 \\
 0 & i
\end{pmatrix*}
\begin{pmatrix*}[r]
 0&-1\\
 1& 0
\end{pmatrix*}
=
\begin{pmatrix*}[r]
0&i\\
i&0
\end{pmatrix*}
=
-\boldsymbol{i},
\notag
\\
\boldsymbol{k}\boldsymbol{i}
&=
\begin{pmatrix*}[r]
-i & 0 \\
 0 & i
\end{pmatrix*}
\begin{pmatrix*}[r]
 0&-i\\
-i& 0
\end{pmatrix*}
=
\begin{pmatrix*}[r]
0 & -1 \\
1 & 0
\end{pmatrix*}
=
\boldsymbol{j},
&
\boldsymbol{i}\boldsymbol{k}
&=
\begin{pmatrix*}[r]
 0&-i\\
-i& 0
\end{pmatrix*}
\begin{pmatrix*}[r]
-i & 0 \\
 0 & i
\end{pmatrix*}
=
\begin{pmatrix*}[r]
 0 & 1 \\
-1 & 0
\end{pmatrix*}
=
-\boldsymbol{j}.
\notag
\end{align}
Wir zu erwarten ist, sind die Produkte der Matrizen
$\boldsymbol{i}$,
$\boldsymbol{j}$
und
$\boldsymbol{k}$
nicht kommutativ.
Diese Relationen können auch in die Gleichungen
\begin{equation}
\boldsymbol{i}^2
=
\boldsymbol{j}^2
=
\boldsymbol{k}^2
=
\boldsymbol{i}
\boldsymbol{j}
\boldsymbol{k}
=
-1
\label{buch:komplex:zahlen:eqn:hamiltonrelationen}
\end{equation}
zusammengefasst werden.
Multipliziert man die letzte Identität von rechts mit $\boldsymbol{k}$,
erhält man
\input{chapters/010-komplex/fig/fig-hamilton.tex}%
\begin{equation*}
-\boldsymbol{k}
=
(
\boldsymbol{i}
\boldsymbol{j}
\boldsymbol{k}
)
\boldsymbol{k}
=
\boldsymbol{i}
\boldsymbol{j}
\boldsymbol{k}^2
=
-
\boldsymbol{i}
\boldsymbol{j}
\quad\Rightarrow\quad
\boldsymbol{i}
\boldsymbol{j}
=
\boldsymbol{k},
\end{equation*}
also die erste Identität in \eqref{buch:komplex:zahlen:eqn:quatrel}.
Die anderen Identitäten können auf ähnliche Weise aus
\eqref{buch:komplex:zahlen:eqn:hamiltonrelationen}
gewonnen werden.

Die Menge
\[
\mathbb{H}
=
\langle 1,
\boldsymbol{i},
\boldsymbol{j},
\boldsymbol{k}
\rangle
=
\{
x_0
+
x_1\boldsymbol{i}
+
x_2\boldsymbol{j}
+
x_3\boldsymbol{k}
\mid
x_0,\dots,x_3\in\mathbb{R}
\}
\subset
M_{2\times 2}(\mathbb{C})
\]
ist einerseits ein reeller Vektorraum und andererseits eine
Algebra mit den Rechenregeln
\eqref{buch:komplex:zahlen:eqn:hamiltonrelationen},
die William Rowan Hamilton im Jahr 1843 entdeckt hat.
Sie heisst die Algebra der \emph{Quaternionen}.
Hamilton sollen die algebraischen Regeln
\index{Hamilton, William Rowan}%
\eqref{buch:komplex:zahlen:eqn:hamiltonrelationen}
während eines Spaziergangs über die Broom-Bridge in Dublin
\index{Dublin}%
\index{Broom-Bridge}%
eingefallen sein und er soll sie spontan in einen Stein
geritzt haben, woran die in Abbildung~\ref{buch:komplex:fig:hamilton}
abgebildete, an der Broom-Bridge angebrachte Plakette erinnert.

Die Matrizen $\boldsymbol{i}$ und $\boldsymbol{k}$ enthalten
Einträge, die selbst komplexe Zahlen sind.
Um eine Darstellung ganz ohne komplexe Zahlen zu bekommen, muss
jeder Eintrag der Matrizen in \eqref{buch:komplex:zahlen:eqn:quaternionen}
durch eine $2\times 2$-Matrix ersetzt werden.
Dadurch entstehen drei $4\times 4$-Matrizen, die wir mit grossen
fetten Buchstaben schreiben:
\begin{align*}
\boldsymbol{I}
&=
\begin{pmatrix*}[r]
 0& 0& 0& 1\\
 0& 0&-1& 0\\
 0& 1& 0& 0\\
-1& 0& 0& 0
\end{pmatrix*},
&
\boldsymbol{J}
&=
\begin{pmatrix*}[r]
 0& 0&-1& 0\\
 0& 0& 0&-1\\
 1& 0& 0& 0\\
 0& 1& 0& 0
\end{pmatrix*}
&&\text{und}
&
\boldsymbol{K}
&=
\begin{pmatrix*}[r]
 0& 1& 0& 0\\
-1& 0& 0& 0\\
 0& 0& 0&-1\\
 0& 0& 1& 0
\end{pmatrix*}.
\end{align*}
Man kann nachrechnen, dass die für
$\boldsymbol{i}$,
$\boldsymbol{j}$
und
$\boldsymbol{k}$
gefundenen Rechenregeln
\eqref{buch:komplex:zahlen:eqn:hamiltonrelationen}
auch für
$\boldsymbol{I}$,
$\boldsymbol{J}$
und
$\boldsymbol{K}$
gelten.

So wie Drehungen der komplexen Ebene durch Multiplikation mit
komplexen Zahlen beschrieben werden, können Quaternionen Drehungen
des dreidmensionalen Raumes beschreiben.
Dies wird im Kapitel~\ref{chapter:qa} beschrieben.




